% \section{Approach: Modeling Transformer Learning Dynamics through an Evolutionary Lens}
% \section{Learning dynamics through an evolutionary lens}
\section{Environmental predictability}

Our general approach relies on drawing a correspondence between aspects of biological adaptation and the learning dynamics observed in Transformers. Fundamentally, we view both systems as optimizing for environmental predictability: the strategy selected depends on the timescale at which information is most reliable. To make these parallels concrete, we consider two examples that illustrate the specific phenomena we investigate.

% Our general approach relies on drawing a correspondence between aspects of biological adaptation and the learning dynamics observed in Transformers. Just as ICL and its underlying dynamics are emergent products of IWL, biological plasticity is an emergent property of evolution. To make these parallels concrete, we consider two classic biological examples that illustrate the specific phenomena we investigate in this work.

\textbf{Plasticity.} The classic example of plasticity is the defense adaptation of the water flea, Daphnia \cite{tollrian1999ecology}. When Daphnia detect chemical cues (kairomones) released by predators, they develop a protective helmet and tail spine. In the absence of these cues, individuals with the exact same genotype develop a standard, non-defensive morphology to conserve energy.
Just as the Daphnia phenotype is determined by the interaction of a fixed genotype with an environmental cue, a Transformer’s output is determined by the interaction of fixed weights with the context provided in the prompt. In both cases, the system adapts its behavior to the current environment without modifying its underlying genotype or weights.

\textbf{Genetic assimilation.} The phenomenon where ICL emerges early in training but is later superseded by IWL parallels genetic assimilation, a concept famously demonstrated by Waddington’s selection experiments on Drosophila \cite{waddington1953genetic}. In these experiments, Waddington exposed fruit fly pupae to a heat shock (an environmental cue), which caused a fraction of them to develop a cross-veinless wing structure (a plastic response). By selectively breeding only those individuals that exhibited the trait, he eventually produced a lineage where the cross-veinless trait appeared without the heat shock.
Similarly, a Transformer initially reliant on prompts (ICL) will, given stable training conditions, consolidate the capability into its weights (IWL), rendering the inference insensitive to the contextual cues that originally guided it \cite{singh2023transient}.

At its core, these two phenomena hinge on distinct aspects of environmental predictability: the reliability of cues to guide adaptation within a lifetime, and the stability of the environment to facilitate evolution across generations. Both dimensions ultimately derive from the predictive power of information at different timescales --- specifically, how accurately experience (whether gathered within a single lifetime or accumulated across generations) anticipates optimal behavior. If this correspondence between biological adaptation and learning in Transformers holds, then systematically manipulating predictability at different timescales (e.g., via noise) should allow us to determine the preferred learning strategy in Transformers. The following section details how we operationalize this experimentally in two controlled training environments.

% Our approach leverages analogies between biological adaptation and Transformer learning to investigate the influence of environmental predictability. The core idea is that the emergence, persistence, and transitions between ICL and IWL in Transformers can be understood as parallels to how organisms evolve to use phenotypic plasticity versus genetically determined traits, and how these strategies themselves can shift (e.g., via genetic assimilation).

% To make these parallels concrete:

% \textbf{ICL as Phenotypic Plasticity:} Consider the water flea (Daphnia), which develops protective helmets and spines (phenotype) only when detecting chemical cues from predators (environmental cue) \cite{tollrian1999ecology}. Similarly, many animals adjust coat thickness and color based on changing daylight hours -- a cue predictive of seasonal temperature and snowfall \cite{underwood_photoperiod_1980, otte_snow_2024}. This rapid, reversible, cue-driven adaptation mirrors how a Transformer uses in-context examples (cue) to tailor its output (phenotype) to a novel task without altering its weights (genotype).

% \textbf{IWL as Genetically Determined/Canalized Traits:} The innate, species-specific song of certain songbirds develops reliably, even in acoustic isolation. This exemplifies a genetically underpinned trait that is invariant to environmental variation, a phenomenon known as environmental canalization \cite{marler2014instinct}. Analogously, IWL captures fundamental, consistently useful knowledge (e.g., grammar for an LLM) that is ``hardcoded'' into model weights through sustained exposure to stable patterns.

% \textbf{ICL Transience as Genetic Assimilation:} Waddington's classic experiments showed that a heat-shock-induced crossveinless wing phenotype in Drosophila (a plastic response), after generations of selection, became genetically assimilated, appearing even without the heat shock \cite{waddington1953genetic}. Analogously, if a Transformer initially uses ICL for a task type that later becomes a stable and frequent part of training, the model might ``assimilate'' this solution into its weights (IWL), potentially reducing reliance on ICL for that specific pattern.

% Guided by this framework, particularly the established role of environmental stability and cue reliability in shaping biological adaptation, our experimental design systematically manipulates these two factors within the training regime of Transformers.
% For cue reliability, we vary the clarity or consistency of the in-context examples provided within each prompt.
% % For instance, in the Omniglot classification task, this involves manipulating the consistency of label-to-image pairings in the few-shot examples. In the sinusoid regression task, it could involve varying the noise or number of example points defining the target function within a prompt.
% For environmental stability, we manipulate the consistency of the underlying task rule across successive training batches. 
% % In the classification task, a stable environment might always present the same underlying (but initially unknown) mapping of image classes to target labels across many prompts, while an unstable environment would change this mapping frequently. For regression, stability refers to how consistently the parameters of the target sinusoid function (e.g., phase, amplitude) are drawn from the same distribution or remain fixed across prompts.

% \section{Background and Related Work}

% \stephanie{Given that we need to save space, I think we could cut this  Related Work section, and move the most important refs/content to the intro. It is partially redundant with the intro, as is.}

% \subsection{In-Context Learning vs. In-Weights Learning}

% \stephanie{we might need to cut this subsection completely}

% In-context learning (ICL) is a key emergent capability of Transformers \cite{vaswani2017attention}, referring to their ability to adapt to novel tasks specified by examples within an input prompt, without requiring gradient-based parameter updates \cite{brown2020language, dong2022survey, liu2023pre}. The mechanisms proposed to underpin ICL are diverse. Prominent theories include the role of induction heads, which copy and complete patterns from the context \cite{olsson2022context, elhage2021mathematical, akyurek2024higher, wang2024Transformers}, and the more recently identified function vector (FV) heads, which learn to compute and apply compact task representations \cite{todd2023function, hendel2023context, le2024disentangling}. Other interpretations frame ICL as an implicit optimization process (e.g., akin to gradient descent) performed within the forward pass \cite{von2023Transformers, dai2022can, akyurek2022learning, ding2023Transformer}, or as a form of implicit Bayesian inference \cite{xie2021explanation, wang2023Transformers, prystawski2023think, zhang2024trained, li2023Transformers}. The training regimen that gives rise to ICL is often conceptualized as a meta-learning process, effectively learning an inference-time learning algorithm \cite{finkbeiner2024chip, kirsch2022general, melo2022Transformers, wang2024Transformers}.

% Conversely, in-weights learning (IWL) describes the conventional knowledge acquisition process in neural networks. Here, task-relevant information is gradually encoded into the model's parameters via gradient-based optimization during training \cite{sun_learning_2024, ruis_procedural_2024, zucchet_how_2025}. IWL typically yields more stable and generalized knowledge for tasks that are frequently encountered or possess a consistent structure. A central aim of our research is to understand how the balance and developmental trajectory of these two learning modes are shaped by the statistical properties of the task environment encountered during training.

% \subsection{Learning Dynamics: Transience, Competition, and Data Influence}

% The preference for ICL vs. IWL is notably dynamic, evolving throughout the training process. A significant observation from recent literature is the phenomenon of ICL transience: ICL capabilities often emerge relatively early in training but may subsequently diminish or be superseded by IWL strategies \cite{singh2023transient}. This shift is particularly evident when tasks are encountered repeatedly or become highly predictable within the training data \cite{anand2024dual, he2024context, park2024understanding}.

% The emergence and relative balance of ICL versus IWL are also influenced by broader characteristics of the data distribution. Factors such as the burstiness of token appearance and Zipfian class distributions, which are typical of natural language, have been shown to affect this balance \cite{chan2022data, reddy2024data}. Most relevant to our current study, input noise and dynamic input-label mappings can impact task predictability and also the balance between ICL and IWL \cite{chan2022data}. Studies investigating these aspects have provided compelling initial evidence that variations in predictability can shift the preference between ICL and IWL.

% Mechanistically, the transition between ICL and IWL is hypothesized to involve both competition and cooperation among neural circuits within the Transformer architecture \cite{singh2023transient, park2024understanding, singh_strategy_2025}. For instance, IWL might predominantly engage parameters within MLP layers, whereas ICL often leverages attention-based circuits, such as induction or FV heads \cite{olsson2022context,todd2023function}. While these distinct pathways might compete for representational capacity or influence, IWL-based structures can also scaffold or facilitate the development of ICL mechanisms.

% \stephanie{(optional) add a small section on predictability in Transformers. E.g. Elmoznino 2024, Goldblum 2024} \alex{I think we should put this in the part in future directions where we suggest that "relative-cost" can be characterized in many ways, including complexity.}

% \subsection{Biological Adaptation: Phenotypic Plasticity and Environmental Predictability}

% Biological organisms have evolved diverse strategies to contend with the challenges posed by variable and uncertain environments. A cornerstone of this adaptability is \emph{phenotypic plasticity}: the capacity of a single genotype to produce a spectrum of different phenotypes—spanning morphology, physiology, and behavior—in direct response to specific environmental cues encountered during an individual's development or lifetime \cite{west2003developmental, pigliucci2001phenotypic, scheiner1993genetics}. This facility allows individuals to tailor their characteristics to local conditions, potentially enhancing survival and reproduction relative to organisms with rigidly fixed phenotypes, particularly in heterogeneous or fluctuating settings. Learning, wherein individuals modify their behavior based on experience to navigate and exploit their surroundings, represents a particularly sophisticated form of such behavioral plasticity \cite{dukas1998evolutionary, snell2019evolution}. Phenotypic plasticity is often contrasted with \emph{environmental canalization}, which is the invariance of phenotypes to environmental variation \cite{waddington_canalization_1942, siegal_waddingtons_2002}.

% The existence and evolutionary trajectory of phenotypic plasticity are intricately linked to several dimensions of environmental predictability:

% \textbf{Environmental Stability.} The degree of environmental variation across generations profoundly influences whether plasticity is favored over more canalized (fixed) strategies \cite{boyce2006predicting, leung2020evolution}. In highly stable environments, where the optimal phenotype seldom changes, natural selection typically favors the evolution of a single, well-adapted fixed phenotype. Maintaining the biological machinery for plasticity in such constant scenarios would likely incur unnecessary costs \cite{dewitt1998costs}.

% \textbf{Cue Reliability.} For plasticity to be adaptive, environmental cues must reliably forecast the selective pressures an organism will face when the corresponding phenotype is expressed \cite{dall2015learning, lehtonen2021sensitive, scheiner2021loss}. If cues are noisy or uninformative, plastic responses may be maladaptive.

% A crucial long-term dynamic in this context is \emph{genetic assimilation}. If an environment that previously favored a plastic response becomes stable for an extended duration, the consistently advantageous phenotype, initially induced by cues, can become genetically encoded. This process results in the phenotype being "constitutively expressed" (i.e., activated even in the absence of the original cue), thereby shifting the trait from a flexible, plastic state to a fixed, genetically determined one \cite{waddington1953genetic, west2003developmental, pigliucci2006genetic}.

% \section{Approach: Tying Learning Dynamics to Evolutionary Phenomena}

% \stephanie{this is also a bit redundant with the intro, and can probably be cut / incorporated into there if needed}

% The rich theoretical framework of evolutionary biology
% % , particularly concerning phenotypic plasticity and its adaptation to environmental predictability,
% offers a potent explanatory lens for understanding the dynamic interplay between ICL and IWL in Transformers. We propose that drawing analogies between these domains can provide computational-level insights \cite{marr82vision} into why specific learning strategies emerge, persist, or transition within these models.

% Our central thesis posits a direct correspondence between these learning modalities and established biological adaptation strategies. ICL strongly mirrors phenotypic plasticity. A classic biological illustration is the development of defensive structures in water fleas (Daphnia). When these organisms detect chemical cues (kairomones) from predators, individuals can develop protective helmets and elongated tail spines, significantly enhancing their survival prospects without any genetic change. In the absence of such cues, they forgo these energetically costly structures \cite{agrawal1999specific_study_if_available, tollrian1999predator_cues_daphnia}. This rapid, cue-driven, and often reversible adaptation is strikingly analogous to how a Transformer, when provided with a few demonstration examples in its context window (the "cue"), tailors its output (the "phenotype") to a novel task---be it translation, sentiment analysis, or code generation---without undergoing retraining, thereby leaving its "genotype" or weights unchanged.

% Conversely, IWL finds its parallel in genetically determined or canalized traits. An example from the natural world is the innate, species-specific song of many songbirds. While some avian species exhibit significant song learning (a form of plasticity), others possess highly stereotyped songs that develop reliably even when individuals are reared in acoustic isolation from conspecifics \cite{marler1991birdsong_innateness_example; e.g., some flycatchers or other species with highly innate songs}. This points to a strong genetic underpinning and developmental canalization, ensuring the consistent structure of the song, which is often vital for species recognition and successful mating. In a similar fashion, IWL captures fundamental, consistently useful knowledge---such as the grammatical rules of a language for an LLM, or core visual features for an object recognition model---that the model applies broadly, representing a "hardcoded" adaptation resulting from sustained exposure to stable patterns within the training data.

% Furthermore, the phenomenon of ICL transience closely resembles the biological process of genetic assimilation. The foundational experiments by Waddington (1953) with Drosophila provide a clear demonstration: by repeatedly exposing fruit fly pupae to a heat shock, he induced a non-standard, crossveinless wing phenotype (a plastic response). After many generations of selectively breeding only those flies that exhibited this heat-induced response, a strain eventually emerged that developed these crossveinless wings reliably, even without the heat shock stimulus \cite{waddington1953genetic}. The acquired characteristic had become genetically assimilated. Analogously, if a Transformer initially employs ICL to address a particular task type (e.g., a specific format of few-shot reasoning problem), but this task type becomes a highly frequent and stable component of its training regimen, the model may progressively "assimilate" the solution into its weights via IWL. This effectively hardcodes the solution for more efficient or robust recall, potentially diminishing reliance on the ICL mechanism for that specific pattern, even if the explicit in-context cues for it become less prominent or are varied in subsequent encounters.

% Central to the evolution of these diverse biological strategies is the predictability of the environment. In broad terms, phenotypic plasticity is favored by selection when environmental conditions fluctuate over time or space, yet reliable cues are available to forecast the optimal phenotype and guide an appropriate adaptive response \cite{scheiner2021loss, dall2015learning, stamps2020information}. Conversely, highly stable environments, where the optimal phenotype rarely changes from generation to generation, tend to select for canalized (fixed) traits, as maintaining the biological machinery for plasticity would incur costs without conferring significant benefits \cite{dewitt1998costs, leung2020reduced}. Moreover, if an initially variable environment that favored a plastic response transitions to a prolonged period of stability where one particular phenotype is consistently advantageous, genetic assimilation can then drive the fixation of this previously plastic trait.

% In our work, we operationalize analogous dimensions of predictability within the training data of Transformers, specifically distinguishing between cue reliability (the clarity and informativeness of information within a single prompt for specifying the required task) and environmental stability (the consistency or predictability of the task structure across different prompts encountered over time). We then systematically investigate whether these factors influence the balance and developmental dynamics of ICL and IWL in a manner consistent with the predictions derived from evolutionary biology.
